\documentclass{amsbook}

\usepackage{enumitem}
\usepackage{amssymb}
\usepackage{hyperref}
\usepackage{mathtools}
\usepackage{geometry}                % See geometry.pdf to learn the layout options. There are lots.
\geometry{a4paper}                   % ... or a4paper or a5paper or ... 

\newtheorem{theorem}{Theorem}[chapter]
\newtheorem{lemma}[theorem]{Lemma}

\theoremstyle{definition}
\newtheorem{definition}[theorem]{Definition}
\newtheorem{example}[theorem]{Example}
\newtheorem{xca}[theorem]{Exercise}

\theoremstyle{remark}
\newtheorem{remark}[theorem]{Remark}

\numberwithin{section}{chapter}
\numberwithin{equation}{chapter}

\begin{document}

\frontmatter

\title{Minqi's Solutions to the Exercises of Sheldon Axler (2014), Linear Algebra Done Right (3rd Edition)}

\author{Minqi Pan\footnote{Website: http://www.minqi-pan.com}}

\maketitle

\setcounter{page}{4}

\tableofcontents

\chapter*{Preface}

Remarkable exercises are:
\begin{enumerate}
\item\hyperref[Exercise 1.C.9]{Exercise 1.C.9: Periodic Functions is NOT a Subspace of $\mathbb{R}^\mathbb{R}$}.
\item\hyperref[Exercise 1.C.13]{Exercise 1.C.13: When the Union of 3 Subspaces is a Subspace}.
\item\hyperref[Exercise 2.B.7]{Exercise 2.B.7: Constructing a Counterexample}.
\item\hyperref[Exercise 2.C.9]{Exercise 2.C.9: $\dim\text{span}(v_1+w,\dots,v_m+w)\geqslant m-1$}.
\item\hyperref[Exercise 3.C.5]{Exercise 3.C.5: Find Bases for $V$ from $\mathcal{M}(T)$ where $T\in\mathcal{L}(V,W)$}.
\item\hyperref[Exercise 3.D.4]{Exercise 3.D.4: $\text{null }T_1=\text{null }T_2$ if and only if $T_1=ST_2$}.
\item\hyperref[Exercise 3.D.5]{Exercise 3.D.5: $\text{range }T_1=\text{range }T_2$ if and only if $T_1=T_2S$}.
\end{enumerate}

\medskip
\begin{flushright}
Minqi Pan\\
\today
\end{flushright}

\mainmatter

\include{minqi_s_solutions_to_the_exercises_of_sheldon_axler_2014_linear_algebra_done_right_3rd_edition_ch1}
\include{minqi_s_solutions_to_the_exercises_of_sheldon_axler_2014_linear_algebra_done_right_3rd_edition_ch2}
\chapter{The Vector Space of Linear Maps}
\section{Exercise 3.A.1: Conditions for a Mapping to be Linear}
Suppose $b,c\in\mathbb{R}$. Define $T:\mathbb{R}^3\to\mathbb{R}^2$ by\[
T(x,y,z)=(2x-4y+3z+b,6x+cxyz).
\]Show that $T$ is linear if and only if $b=c=0$.
\begin{proof}
\begin{enumerate}
\item Suppose $b=c=0$, then\[
T(x,y,z)=(2x-4y+3z,6x).
\]

Pick $u=(u_1,u_2,u_3),v=(v_1,v_2,v_3)$, then\[
\begin{split}
T(u+v)&=(2(u_1+v_1)-4(u_2+v_2)+3(u_3+v_3),6(u_1+v_1))\\
&=(2u_1+2v_1-4u_2-4v_2+3u_3+3v_3,6u_1+6v_1)
\end{split}
\]and\[
\begin{split}
Tu+Tv&=(2u_1-4u_2+3u_3,6u_1)+(2v_1-4v_2+3v_3,6v_1)\\
&=(2u_1-4u_2+3u_3+2v_1-4v_2+3v_3,6u_1+6v_1).
\end{split}
\]Hence $T(u+v)=Tu+Tv$.

Pick $\lambda\in\mathbb{F}$, then\[
\begin{split}
T(\lambda v)&=T(\lambda u_1,\lambda u_2,\lambda u_3)\\
&=(2\lambda u_1-4\lambda u_2+3\lambda u_3,6\lambda u_1)
\end{split}
\]and\[
\begin{split}
\lambda T(v)&=\lambda(2u_1-4u_2+3u_3,6u_1)\\
&=(2\lambda u_1-4\lambda u_2+3\lambda u_3,6\lambda u_1).
\end{split}
\]Hence $T(\lambda v)=\lambda T(v)$.

\item By definition, if $T$ is linear, then $T(u+v)=Tu+Tv$ for all $u,v\in\mathbb{R}^3$. In particular, take $u=(1,1,1),v=(1,1,1)$. Then\[
\begin{split}
T(u+v)&=(2\cdot 2-4\cdot 2+3\cdot 2+b,6\cdot 2+c\cdot 2\cdot 2\cdot 2)\\
&=(2+b,12+8c)\\
&=Tu+Tv\\
&=(2-4+3+b+2-4+3+b,6+c+6+c)\\
&=(2+2b,12+2c)
\end{split}
\]Hence $b=2b$ and $8c=2c$, we thus conclude that $b=c=0$.
\end{enumerate}
\end{proof}
\section{Exercise 3.A.5: $\mathcal{L}(V, W)$ is a vector space}
Prove the assertion in 3.7 of \cite{axler2014linear}: $\mathcal{L}(V, W)$ is a vector space.
\begin{proof}
Define $\mathbf{0}$ as the linear map in $\mathcal{L}(V,W)$ such that\[\
\mathbf{0}(v)=0\quad\forall v\in V.
\]

Pick $R,S,T\in\mathcal{L}(V,W), v\in V$ and $a,b\in\mathbb{F}$.

\begin{enumerate}
\item Commutativity: The commutativity of vector space $W$ implies $Sv+Tv=Tv+Sv$. Hence $S+T=T+S$.
\item Associativity: The associativity of vector space $W$ implies\[(Rv+Sv)+Tv=Rv+(Sv+Tv)\]and\[(ab)Sv=a(bSv).\] Hence $(R+S)+T=R+(S+T)$ and $(ab)S=a(bS)$.
\item Additive identity: It is clear that $T+\mathbf{0}=T$.
\item Additive inverse: Define $-T$ as the linear map in $\mathcal{L}(V,W)$ such that\[
-T(v)=-(Tv)\quad\forall v\in V.
\]

Then $T+(-T)=\mathbf{0}$.
\item Multiplicative identity: The multiplicative identity of vector space $W$ implies $(1T)(v)=1(Tv)=Tv$. Hence $1T=T$.
\item Distributive properties: The distributive properties of $W$ implies\[a(Sv+Tv)=aSv+aTv\]and\[(a+b)Sv=aSv+bSv.\]Hence $a(S+T)=aS+aT$ and $(a+b)S=aS+bS$.
\end{enumerate}
\end{proof}
\section{Exercise 3.A.6: Algebraic Properties of Products of Linear Maps}
Prove the assertions in 3.9 of \cite{axler2014linear}: The associativity, identity and distributive properties of Products of Linear Maps.
\begin{proof}
\begin{enumerate}
\item Associativity: pick linear maps $T_1, T_2$ and $T_3$ such that $T_3$ maps into the domain of $T_2$, and $T_2$ maps into the domain of $T_1$. Then for any $v$ in the domain of $T_3$,\[
((T_1T_2)T_3)(v)=(T_1T_2)(T_3(v))=T_1(T_2(T_3(v)))=T_1(T_2T_3)(v).
\]Hence\[
(T_1T_2)T_3=T_1(T_2T_3).
\]
\item Identity: Pick $T\in\mathcal{L}(V,W)$. Define $I_V$ as the identity map on $V$ and $I_W$ as the identity map on $W$. Then for any $v\in V$,\[
(TI_V)(v)=T(I_V(v))=T(v)=I_W(T(v))=(I_WT)(v).
\]Hence\[
TI_V=I_WT=T.
\]
\item Distributive properties. Pick $T,T_1,T_2\in\mathcal{L}(U,V)$ and $S,S_1,S_2\in\mathcal{L}(V,W)$. Then for any $u\in U$,\[
\begin{split}
((S_1+S_2)T)(u)&=S_1(T(u))+S_2(T(u))\\
&=(S_1T)(u)+(S_2T)(u)\\
&=(S_1T+S_2T)(u)
\end{split}
\]and\[
\begin{split}
(S(T_1+T_2))(u)&=S(T_1(u)+T_2(u))\\
&=S(T_1(u))+S(T_2(u))\\
&=(ST_1)(u)+(ST_2)(u)\\
&=(ST_1+ST_2)(u).
\end{split}
\]Hence\[
(S_1+S_2)T=S_1T+S_2T
\]and\[
S(T_1+T_2)=ST_1+ST_2.
\]
\end{enumerate}
\end{proof}
\section{Exercise 3.A.7: Linear Map on 1D Space is Scalar Multiplication}
Show that every linear map from a $1$-dimensional vector space to itself is multiplication by some scalar. More precisely, prove that if $\dim V=1$ and $T\in\mathcal{L}(V,V)$, then there exists $\lambda\in\mathbb{F}$ such that $Tv=\lambda v$ for all $v\in V$.
\begin{proof}
Since $\dim V=1$, there exists a vector $b$ which is a basis of $V$. Pick any $T\in\mathcal{L}(V,V)$, then, by the definition of basis, there exists $\lambda\in\mathbb{F}$ such that\[
T(b)=\lambda b.
\]We claim that $Tv=\lambda v$ for all $v\in V$. To show this, pick any $v\in V$, then there exists $\alpha\in\mathbb{F}$ such that\[
v=\alpha b.
\]Hence\[
T(v)=T(\alpha b)=\alpha T(b)=\alpha \lambda b=\lambda\alpha b=\lambda v.
\]
\end{proof}

\section{Exercise 3.A.8: Homogeneity Alone does NOT Imply Linearity}
Give an example of a function $\varphi:\mathbb{R}^2\to\mathbb{R}$ such that\[
\varphi(av)=a\varphi(v)
\]for all $a\in\mathbb{R}$ and all $v\in\mathbb{R}^2$ but $\varphi$ is not linear.
\begin{proof}
Let $\varphi(x,y)\equiv(x^3+y^3)^\frac{1}{3}$ where $(x,y)\in\mathbb{R}^2$.
\begin{enumerate}
\item Pick any $a\in\mathbb{R}$, then\[
\begin{split}
\varphi(ax,ay)&=(a^3x^3+a^3y^3)^\frac{1}{3}\\
&=(a^3(x^3+y^3))^\frac{1}{3}\\
&=a(x^3+y^3)^\frac{1}{3}\\
&=a\varphi(x,y).
\end{split}
\]
\item $\varphi$ is not linear because\[
\begin{split}
\varphi((1,2)+(3,4))&=(4^3+6^3)^\frac{1}{3}\\
&=280^\frac{1}{3}\\
&\ne9^\frac{1}{3}+91^\frac{1}{3}\\
&=\varphi(1,2)+\varphi(3,4)
\end{split}
\]
\end{enumerate}
\end{proof}

\section{Exercise 3.A.9: Additivity Alone does NOT Imply Linearity}
Give an example of a function $\varphi:\mathbb{C}\to\mathbb{C}$ such that\[
\varphi(w+z)=\varphi(w)+\varphi(z)
\]for all $w,z\in\mathbb{C}$ but $\varphi$ is not linear.
\begin{proof}
Let $\varphi(z)=\Re(z)$ where $z\in\mathbb{C}$.
\begin{enumerate}
\item Pick any $w,z\in\mathbb{C}$, then\[
\begin{split}
\varphi(w+z)&=\Re(w+z)\\
&=\Re(w)+\Re(z)\\
&=\varphi(w)+\varphi(z).
\end{split}
\]
\item $\varphi$ is not linear because\[
\begin{split}
i\varphi(i)&=i\Re(i)\\
&=0\\
&\ne -1\\
&=\Re(-1)\\
&=\varphi(i\cdot i).
\end{split}
\]
\end{enumerate}
\end{proof}

\chapter{Null Spaces and Ranges}
\section{Exercise 3.B.1: A Linear Map $T$ with $\dim\text{null } T=3,\dim\text{range } T=2$}
Give an example of a linear map $T$ such that \[
\dim\text{null } T=3\text{ and }\dim\text{range } T=2.
\]
\begin{proof}
Define $T$ as the linear map in $\mathcal{L}(\mathbb{R}^5,\mathbb{R}^5)$ such that\[
T(a,b,c,d,e)=(a,b,0,0,0)
\]for all $(a,b,c,d,e)\in\mathbb{R}^5$.

It is clear that $\text{range }T=2$ and $\dim V=5$. Hence, by the fundamental theorem of linear maps,\[
\dim\text{null }T=\dim V-\dim\text{range }T=3.
\]
\end{proof}
\section{Exercise 3.B.2: $\text{range }S\subset\text{null }T$ Implies $(ST)^2=0$}
Suppose $V$ is a vector space and $S,T\in\mathcal{L}(V,V)$ are such that\[
\text{range }S\subset\text{null }T.
\]Prove that $(ST)^2=0$.
\begin{proof}
Pick any $v\in V$, then
\begin{gather*}
v\in V\\
\ArrowBetweenLines[\Downarrow]
T(v)\in V\\
\ArrowBetweenLines[\Downarrow]
S(T(v))\in\text{range } S\\
\ArrowBetweenLines[\Downarrow]
S(T(v))\in\text{null }T\\
\ArrowBetweenLines[\Downarrow]
T(S(T(v)))=0\\
\ArrowBetweenLines[\Downarrow]
S(T(S(T(v))))=0.
\end{gather*}
Hence $(ST)^2=0$.
\end{proof}
\section{Exercise 3.B.3: Spanning = Surjective; Independent = Injective}
Suppose $v_1,\dots,v_m$ is a list of vectors in $V$. Define $T\in\mathcal{L}(\mathbb{F}^m,V)$ by\[
T(z_1,\dots,z_m)=z_1v_1+\dots+z_mv_m.
\]
\begin{enumerate}[label=(\alph*)]
\item What property of $T$ corresponds to $v_1,\dots,v_m$ spanning $V$?
\item What property of $T$ corresponds to $v_1,\dots,v_m$ being linearly independent?
\end{enumerate}
\begin{proof}
\begin{enumerate}[label=(\alph*)]
\item We claim that $T$ being surjective corresponds to $v_1,\dots,v_m$ spanning $V$.

If $T$ is surjective, then $\text{range }T=V$. Since $\text{range }T=\text{span}(v_1,\dots,v_m)$, we conclude that $v_1,\dots,v_m$ spans $V$.

If $v_1,\dots,v_m$ spans $V$, then $\text{range }T=V$. Hence $T$ is surjective.
\item We claim that $T$ being injective corresponds to $v_1,\dots,v_m$ being linearly independent.

If $T$ is injective, then 3.16 of \cite{axler2014linear} implies that $\text{null }T=\{0\}$. Hence $T(z_1,\dots,z_m)=z_1v_1+\dots+z_mv_m=0$ implies $z_1=\dots=z_m=0$. Therefore $v_1,\dots,v_m$ is linearly independent.

If $v_1,\dots,v_m$ is linearly independent, then $z_1v_1+\dots+z_mv_m=0$ implies $z_1=\dots=z_m=0$. Hence $T(z_1,\dots,z_m)=0$ implies $z_1=\dots=z_m=0$. By 3.16 of \cite{axler2014linear}, $T$ is therefore injective.
\end{enumerate}
\end{proof}
\section{Exercise 3.B.4: $\{T\in\mathcal{L}(\mathbb{R}^5,\mathbb{R}^4):\dim\text{null }T>2\}$ is NOT a Subspace}
Show that\[
\{T\in\mathcal{L}(\mathbb{R}^5,\mathbb{R}^4):\dim\text{null }T>2\}
\]
is not a subspace of $\mathcal{L}(\mathbb{R}^5,\mathbb{R}^4)$.
\begin{proof}
Let $S,T$ be the linear maps in $\mathcal{L}(\mathbb{R}^5,\mathbb{R}^4)$ such that\[
S(x_1,x_2,x_3,x_4,x_5)=(x_1,x_2,0,0)
\]and\[
T(x_1,x_2,x_3,x_4,x_5)=(x_1,0,x_3,0)
\]for all $(x_1,x_2,x_3,x_4,x_5)\in\mathbb{R}^5$. Then\[
(S+T)(x_1,x_2,x_3,x_4,x_5)=(2x_1,x_2,x_3,0)
\]for all $(x_1,x_2,x_3,x_4,x_5)\in\mathbb{R}^5$. 

Clearly, $\dim\text{range }S=\dim\text{range }T=2$ and $\dim\text{range }(S+T)=3$. Hence the fundamental theorem of linear maps implies\[
\dim\text{null }S=\dim\text{null }T=5-2=3>2
\]and\[
\dim\text{null }(S+T)=5-3=2\not>2.
\]We have thus shown that\[
\{T\in\mathcal{L}(\mathbb{R}^5,\mathbb{R}^4):\dim\text{null }T>2\}
\] is not closed under addition, hence it is not a subspace of $\mathcal{L}(\mathbb{R}^5,\mathbb{R}^4)$.
\end{proof}
\section{Exercise 3.B.5: A Linear Map $T$ such that $\text{range }T=\text{null }T$}
Give an example of a linear map $T:\mathbb{R}^4\to\mathbb{R}^4$ such that\[
\text{range }T=\text{null }T.
\]
\begin{proof}
Let $T$ be the linear map in $\mathcal{L}(\mathbb{R}^4,\mathbb{R}^4)$ such that\[
T(x_1,x_2,x_3,x_4)=(0, 0, x_1,x_2)
\]for all $(x_1,x_2,x_3,x_4)\in\mathbb{R}^4$. Then\[
\text{range }T=\{(0, 0, x_1,x_2):x_1,x_2\in\mathbb{R}\}
\]and\[
\text{null }T=\{(0,0,x_3,x_4):x_3,x_4\in\mathbb{R}\}.
\]Hence\[
\text{range }T=\text{null }T.
\]
\end{proof}

\chapter{Matrices}
\section{Exercise 3.C.1: $\mathcal{M}(T)$ Has At Least $\dim\text{range }T$ Nonzero Entries}
Suppose $V$ and $W$ are finite-dimensional and $T\in\mathcal{L}(V,W)$. Show that with respect to each choice of bases of $V$ and $W$, the matrix of $T$ has at least $\dim\text{range }T$ nonzero entries.
\begin{proof}
Choose any basis $v_1,\dots,v_n$ of $V$ and basis $w_1,\dots,w_m$ of $W$, suppose the matrix of $T$  with respect to those bases has $k$ nonzero entries where
\begin{equation}
\label{eq3}
k<\dim\text{range }T.
\end{equation}

Then the matrix has at most $k$ nonzero columns. Consequently, by the definition of matrix of a linear map, at most $k$ vectors of $Tv_1,\dots,Tv_n$ are nonzero. Hence\[\dim\text{range }T\leqslant k,\]which contradicts with (\ref{eq3}).
\end{proof}
\section{Exercise 3.C.2: Find Bases from the Matrix of $D\in\mathcal{L}(\mathcal{P}_3(\mathbb{R}),\mathcal{P}_2(\mathbb{R}))$}
Suppose $D\in\mathcal{L}(\mathcal{P}_3(\mathbb{R}),\mathcal{P}_2(\mathbb{R}))$ is the differentiation map defined by $Dp=p'$. Find a basis of $\mathcal{P}_3(\mathbb{R})$ and a basis of $\mathcal{P}_2(\mathbb{R}))$ such that the matrix of $D$ with respect to these bases is\[
\begin{bmatrix}
1&0&0&0\\
0&1&0&0\\
0&0&1&0
\end{bmatrix}.
\]
\begin{proof}
Define\[
w_1\equiv1, w_2\equiv x, w_3\equiv x^2
\]and\[
v_1\equiv x, v_2\equiv\frac{1}{2}x^2, v_3\equiv\frac{1}{3}x^3, v_4\equiv1.
\]

Clearly, $w_1,w_2,w_3$ is a basis of $\mathcal{P}_2(\mathbb{R})$ and $v_1,v_2,v_3,v_4$ is a basis of $\mathcal{P}_3(\mathbb{R})$. Moreover, \[
Dv_1=w_1,Dv_2=w_2,Dv_3=w_3,Dv_4=0.
\]Therefore the matrix of $D$ under these bases is \[
\begin{bmatrix}
1&0&0&0\\
0&1&0&0\\
0&0&1&0
\end{bmatrix}.
\]
\end{proof}
\section{Exercise 3.C.3: Find Bases for $V, W$ from $\mathcal{M}(T)$ where $T\in\mathcal{L}(V,W)$}
Suppose $V$ and $W$ are finite-dimensional and $T\in\mathcal{L}(V,W)$. Prove that there exist a basis of $V$ and a basis of $W$ such that with respect to these bases, all entries of $\mathcal{M}(T)$ are $0$ except that the entries in row $j$, column $j$, equal $1$ for $1\leqslant j\leqslant\dim\text{range }T$.
\begin{proof}
Let $u_1,\dots,u_m$ be a basis of $\text{null }T$; thus $\dim\text{null }T=m$. By 2.33 of \cite{axler2014linear}, the linearly independent list $u_1,\dots,u_m$ can be extended to a basis
\begin{equation}
\label{eq4}
v_1,\dots,v_n,u_1,\dots,u_m
\end{equation}
of $V$. Thus $\dim V=m+n$. It is easy to verify that $Tv_1,\dots,Tv_n$ is a basis of $\text{range }T$. Again, by 2.33 of \cite{axler2014linear}, the linearly independent list $Tv_1,\dots,Tv_n$ can be extended to a basis
\begin{equation}
\label{eq5}
Tv_1,\dots,Tv_n,s_1,\dots,s_l
\end{equation}
of $W$. Thus $\dim W=n+l$. The matrix of $T$ under bases (\ref{eq4}) and (\ref{eq5}) is then of size $(n+l)\times(n+m)$, and all of its entries equal to $0$ except that the entries in row $j$, column $j$, equal $1$ for $1\leqslant j\leqslant n$.
\end{proof}
\section{Exercise 3.C.4: Find Bases for $W$ from $\mathcal{M}(T)$ where $T\in\mathcal{L}(V,W)$}
Suppose $v_1,\dots,v_m$ is a basis of $V$ and $W$ is finite-dimensional. Suppose $T\in\mathcal{L}(V,W)$. Prove that there exists a basis $w_1,\dots,w_n$ of $W$ such that all the entries in the first column of $\mathcal{M}(T)$ (with respect to the bases $v_1,\dots,v_m$ and $w_1,\dots,w_n$) are $0$ except for possibly a $1$ in the first row, first column.
\begin{proof}
If $Tv_1=0$, then any basis of $W$ would make all the entries in the first column of $\mathcal{M}(T)$ equal to $0$.

If $Tv_1\ne 0$, then let $w_1\equiv Tv_1$. By 2.33 of \cite{axler2014linear}, the linearly independent list $w_1$ can be extended to a basis $w_1,\dots,w_n$ of $W$. Under this basis, all the entries in the first column of $\mathcal{M}(T)$ are $0$ except for a $1$ in the first row, first column.
\end{proof}
\section{Exercise 3.C.5: Find Bases for $V$ from $\mathcal{M}(T)$ where $T\in\mathcal{L}(V,W)$}
\label{Exercise 3.C.5}
Suppose $w_1,\dots,w_n$ is a basis of $W$ and $V$ is finite-dimensional. Suppose $T\in\mathcal{L}(V,W)$. Prove that there exists a basis $v_1,\dots,v_m$ of $V$ such that all the entries in the first row of $\mathcal{M}(T)$ (with respect to the bases $v_1,\dots,v_m$ and $w_1,\dots,w_n$) are $0$ except for possibly a $1$ in the first row, first column.
\begin{proof}
Pick any basis\[
v_1,\dots,v_m
\]of $V$. Denote $M$ to be the matrix of $T$ under this basis. If all the entries in the first row of $M$ are $0$, then the proof is done.

Otherwise, suppose column $j$ of the first row of $M$ contains a nonzero entry. Let\[
v_1'\equiv \frac{1}{M_{1,j}}v_j
\]and\[
\{v_2',\dots,v_m'\}\equiv\{v_1,\dots,v_{j-1},v_{j+1},\dots,v_m\}.
\]Then $v_1',\dots,v_m'$ is also a basis of $V$.

Denote $M'$ as the matrix of $T$ under the basis\[
v_1',\dots,v_m'.
\]

Further, let\[
\begin{split}
v_1''&\equiv v_1',\\
v_2''&\equiv v_2'-M'_{1,2}v_1',\\
&\vdots\\
v_m''&\equiv v_m'-M'_{1,m}v_1'.
\end{split}
\]Then $v_1'',\dots,v_m''$ is also a basis of $V$. Moreover,\[
\begin{split}
T(v_1'')=T(v_1')&=T(\frac{1}{M_{1,j}}v_j)\\
&=\frac{1}{M_{1,j}}(M_{1,j}w_1+M_{2,j}w_2+\dots+M_{n,j}w_n)\\
&=w_1+\frac{M_{2,j}}{M_{1,j}}w_2+\dots+\frac{M_{n,j}}{M_{1,j}}w_n.
\end{split}
\]And for any $i=2,\dots,m$,\[
\begin{split}
T(v_i'')&=T(v_i'-M'_{1,i}v_1')\\
&=T(v_i')-M'_{1,i}T(v_1')\\
&=M'_{1,i}w_1+\dots+M'_{n,i}w_n-M'_{1,i}(w_1+\frac{M_{2,j}}{M_{1,j}}w_2+\dots+\frac{M_{n,j}}{M_{1,j}}w_n)\\
&=(M'_{2,i}-\frac{M'_{1,i}M_{2,j}}{M_{1,j}})w_2+\dots+(M'_{n,i}-\frac{M'_{1,i}M_{n,j}}{M_{1,j}})w_n.
\end{split}
\]Hence all the entries in the first row of $\mathcal{M}(T)$ with respect to the bases $v_1'',\dots,v_m''$ and $w_1,\dots,w_n$ are $0$ except for a $1$ in the first row, first column.
\end{proof}

\chapter{Invertibility and Isomorphic Vector Spaces}
\section{Exercise 3.D.1: $(ST)^{-1}=T^{-1}S^{-1}$}
Suppose $T\in\mathcal{L}(U,V)$ and $S\in\mathcal{L}(V,W)$ are both invertible linear maps. Prove that $ST\in\mathcal{L}(U,W)$ is invertible and that $(ST)^{-1}=T^{-1}S^{-1}$.
\begin{proof}
Since $T$ and $S$ are invertible, 3.56 of \cite{axler2014linear} implies that $T$ and $S$ are injective and surjective.

Suppose $S(T(u_1))=S(T(u_2))$ where $u_1,u_2\in U$, then that $S$ is injective implies $T(u_1)=T(u_2)$, and that $T$ is injective implies $u_1=u_2$. Hence $ST$ is injective.

Pick any $w\in W$. That $S$ is surjective implies the existence of a $v\in V$ such that $S(v)=w$. Further, that $T$ is surjective implies the existence of a $u\in U$ such that $T(u)=v$. Hence $S(T(u))=w$. Thus $ST$ is surjective.

Thus 3.56 of \cite{axler2014linear} implies that $ST$ is invertible.

Finally, since\[
(T^{-1}S^{-1})(ST)=T^{-1}(S^{-1}S)T=T^{-1}T=I
\]and\[
(ST)(T^{-1}S^{-1})=S(TT^{-1})S^{-1}=SS^{-1}=I,
\]we conclude that\[
(ST)^{-1}=T^{-1}S^{-1}.
\]
\end{proof}
\section{Exercise 3.D.2: Noninvertible Operators on $V$ is NOT a Subspace}
Suppose $V$ is finite-dimensional and $\dim V>1$. Prove that the set of noninvertible operators on $V$ is not a subspace of $\mathcal{L}(V)$.
\begin{proof}
Let $v_1,v_2,\dots,v_n$ be a basis of $V$ where $n\geqslant 2$. Let $S$ be the operator on $V$ defined as\[
S(a_1v_1+a_2v_2+\dots+a_nv_n)=a_1v_1,
\] and let $T$ be the operator on $V$ defined as\[
T(a_1v_1+a_2v_2+\dots+a_nv_n)=a_2v_2+a_3v_3+\dots+a_nv_n.
\]

Then $S$ is noninvertible because $v_2\in\text{null }S$, and $T$ is noninvertible because $v_1\in\text{null }T$. However, $S+T$ is invertible because $S+T$ equals the identity operator.

We have shown that the set of noninvertible operators on $V$ is not closed under addition, hence it is not a subspace of $\mathcal{L}(V)$.
\end{proof}
\section{Exercise 3.D.3: Injective Subspace Mappings and Invertible Operators}
Suppose $V$ is finite-dimensional, $U$ is a subspace of $V$, and $S\in\mathcal{L}(U,V)$. Prove there exists an invertible operator $T\in\mathcal{L}(V)$ such that $Tu=Su$ for every $u\in U$ if and only if $S$ is injective.
\begin{proof}
Suppose there exists an invertible operator $T\in\mathcal{L}(V)$ such that $Tu=Su$ for every $u\in U$. It follows that
\begin{equation}
\label{eq6}
T^{-1}(Su)=T^{-1}(Tu)=u
\end{equation}
for every $u\in U$. Pick any $u_1,u_2\in U$ such that $Su_1=Su_2$. Then $T^{-1}(Su_1)=T^{-1}(Su_2)$, and (\ref{eq6}) implies that $u_1=u_2$. Hence $S$ is injective.

Suppose $S$ is injective. Pick a basis $u_1,\dots,u_k$ of $U$ and extend it to a basis\[
u_1,\dots,u_k,v_{k+1},\dots,v_n
\]of $V$; thus $\dim U = k, \dim V=n$. Since $S$ is injective, $S$ is an isomorphism between $U$ and $\text{range }S$. Hence \[
Su_1,\dots,Su_k
\]is a basis of $\text{range }T$ and can be extended to a basis\[
Su_1,\dots,Su_k,t_{k+1},\dots,t_{n}
\]of $V$. Define $T$ as the operator on $V$ such that\[
T(u_i)=Su_i\text{ where }i=1,\dots,k
\]and\[
T(v_i)=t_i\text{ where }i=k+1,\dots,n.
\]Clearly, $T$ is invertible and $Tu=Su$ for every $u\in U$.
\end{proof}
\section{Exercise 3.D.4: $\text{null }T_1=\text{null }T_2$ if and only if $T_1=ST_2$}
\label{Exercise 3.D.4}
Suppose $W$ is finite-dimensional and $T_1,T_2\in\mathcal{L}(V,W)$. Prove that $\text{null }T_1=\text{null }T_2$ if and only if there exists an invertible operator $S\in\mathcal{L}(W)$ such that $T_1=ST_2$.
\begin{proof}
Suppose there exists an invertible operator $S\in\mathcal{L}(W)$ such that $T_1=ST_2$, then\[
S^{-1}T_1=T_2.
\]Pick any $v_1\in\text{null }T_1$, then $T_1v_1=0$ and \[
T_2v_1=S^{-1}T_1(v_1)=0.
\]Hence $\text{null }T_1\subset\text{null }T_2$. Also, pick any $v_2\in\text{null }T_2$, then\[
T_1v_2=ST_2(v_2)=0.
\]Hence $\text{null }T_2\subset\text{null }T_1$. Therefore\[
\text{null }T_1=\text{null }T_2.
\]

On the other hand, suppose $\text{null }T_1=\text{null }T_2$. Denote\[
N\equiv\text{null }T_1=\text{null }T_2
\]and\[
m\equiv\dim W.
\]Pick a basis $w_1,\dots,w_k$ of $\text{range }T_2$ and extend it to a basis of $W$:\[
w_1,\dots,w_k,w_{k+1},\dots,w_m.
\]Then there exists nonzero vectors $v_1,\dots,v_k\in V$ such that\[
T_2v_i=w_i,\quad\forall i\in\{1,\dots,k\}.
\]Then it is easy to verify that\[
V=N\oplus\text{span}(v_1,\dots,v_k).
\]and\[
T_1v_1,\dots,T_1v_k
\]is a linearly independent list, which can be extended to a basis of $W$:\[
T_1v_1,\dots,T_1v_k,w_{k+1}',\dots,w_m'.
\]Let $S\in\mathcal{L}(W)$ be defined by\[
S(w_i)\equiv T_1v_i,\quad\forall i\in\{1,\dots,k\},
\]and\[
S(w_i)\equiv w_i',\quad\forall i\in\{k+1,\dots,m\}.
\]We claim that $T_1=ST_2$. To show this, pick any $v\in V$. Then there exists $v_0\in N$ and $a_1,\dots,a_k\in\mathbb{F}$ such that\[
v=v_0+a_1v_1+\dots+a_kv_k
\]and\[
\begin{split}
T_1(v)&=T_1(v_0+a_1v_1+\dots+a_kv_k)\\
&=a_1T_1v_1+\dots+a_kT_1v_k\\
&=a_1Sw_1+\dots+a_kSw_k\\
&=S(a_1w_1+\dots+a_kw_k)\\
&=S(T_2(a_1v_1+\dots+a_kv_k))\\
&=S(T_2(v_0+a_1v_1+\dots+a_kv_k))\\
&=ST_2(v).
\end{split}
\]
\end{proof}
\section{Exercise 3.D.5: $\text{range }T_1=\text{range }T_2$ if and only if $T_1=T_2S$}
\label{Exercise 3.D.5}
Suppose $V$ is finite-dimensional and $T_1,T_2\in\mathcal{L}(V,W)$. Prove that $\text{range }T_1=\text{range }T_2$ if and only if there exists an invertible operator $S\in\mathcal{L}(V)$ such that $T_1=T_2S$.
\begin{proof}
Suppose there exists an invertible operator $S\in\mathcal{L}(V)$ such that $T_1=T_2S$. Pick any $w\in\text{range }T_1$, then there exists $v_1\in V$ such that $T_1v_1=w$. Hence $T_2Sv_1=w$ and $w\in\text{range }T_2$. Also, pick any $w\in\text{range }T_2$, then there exists $v_2\in V$ such that $T_2v_2=w$. Hence $T_1S^{-1}v_2=w$ and $w\in\text{range }T_1$. Therefore $\text{range }T_1=\text{range }T_2$.

On the other hand, suppose $\text{range }T_1=\text{range }T_2$. Denote\[
R\equiv\text{range }T_1=\text{range }T_2
\]and\[
n\equiv\dim V.
\]We claim that $R$ is finite-dimensional. To show this, pick a basis $u_1,\dots,u_k$ of $\text{null }T_1$ and extend it to a basis of $V$:\[
u_1,\dots,u_k,u_{k+1},\dots,u_n.
\]Then it is easy to verify that\[
T_1u_{k+1},\dots,T_1u_n
\]is a basis of $R$; thus\[
\dim R=n-k.
\]The fundamental theorem of linear maps then implies\[
k=\dim\text{null }T_1=\dim V-\dim R=\dim\text{null }T_2.
\]Moreover, there exist $v_{k+1},\dots,v_n\in V$ such that\[
T_2v_i=T_1u_i,\quad\forall i\in\{k+1,\dots,n\}.
\]Pick a basis $v_1,\dots,v_k$ of $\text{null }T_2$. We claim that\[
v_1,\dots,v_k,v_{k+1},\dots,v_n
\]is a basis of $V$. By 2.39 of \cite{axler2014linear}, it suffices to show its linear independence. Suppose there exists $a_1,\dots,a_n\in\mathbb{F}$ such that\[
a_1v_1+\dots+a_nv_n=0.
\]Then
\begin{gather*}
T_2(a_1v_1+\dots+a_nv_n)=0\\
\ArrowBetweenLines[\Downarrow]
a_{k+1}T_2v_{k+1}+\dots+a_nT_2v_n=0\\
\ArrowBetweenLines[\Downarrow]
a_{k+1}T_1u_{k+1}+\dots+a_nT_1u_n=0\\
\ArrowBetweenLines[\Downarrow]
a_{k+1}=\dots=a_n=0\\
\ArrowBetweenLines[\Downarrow]
a_1v_1+\dots+a_kv_k=0\\
\ArrowBetweenLines[\Downarrow]
a_1=\dots=a_k=0.
\end{gather*}
Let $S\in\mathcal{L}(V)$ be defined by\[
S(u_i)\equiv v_i,\quad\forall i\in\{1,\dots,n\}.
\]We claim that $T_1=T_2S$. To show this, pick any $v\in V$. Then there exists $a_1,\dots,a_n\in\mathbb{F}$ such that\[
v=a_1u_1+\dots+a_nu_n.
\]Then\[
\begin{split}
T_1(v)&=T_1(a_1u_1+\dots+a_nu_n)\\
&=a_{k+1}T_1u_{k+1}+\dots+a_nT_1u_n\\
&=a_{k+1}T_2v_{k+1}+\dots+a_nT_2v_n\\
&=T_2(a_{k+1}v_{k+1}+\dots+a_nv_n)\\
&=T_2(a_{k+1}Su_{k+1}+\dots+a_nSu_n)\\
&=T_2(S(a_{k+1}u_{k+1}+\dots+a_nu_n))\\
&=T_2(S(v)).
\end{split}
\]
\end{proof}

\chapter{Products and Quotients of Vector Spaces}
\section{Exercise 3.E.1: $T$ is a Linear Map iff $\text{graph of }T$ is a Subspace of $V\times W$}
Suppose $T$ is a function from $V$ to $W$. The graph of $T$ is the subset of $V\times W$ defined by\[
\text{graph of }T=\{(v,Tv)\in V\times W:v\in V\}.
\]Prove that $T$ is a linear map if and only if the graph of $T$ is a subspace of $V\times W$.
\begin{proof}
Suppose $T$ is a linear map. By 3.11 of \cite{axler2014linear}, $T(0)=0$. Hence\[
(0,0)\in\text{graph of }T.
\]Suppose $(u, Tu), (w,Tw)\in\text{graph of }T$. Then $T(u+w)=Tu+Tw\in W$ implies\[
(u+w,Tu+Tw)\in\text{graph of }T.
\]Suppose $a\in\mathbb{F}$ and $(u, Tu)\in\text{graph of }T$. Then $T(au)=aTu$ implies\[
(au,aTu)\in\text{graph of }T.
\]By 13.4 of \cite{axler2014linear}, we therefore conclude that the graph of $T$ is a subspace of $V\times W$.

On the other hand, suppose the graph of $T$ is a subspace of $V\times W$. Pick any $u,v\in V$ such that $(u,Tu), (v, Tv)\in\text{graph of }T$. Then $(u+v,Tu+Tv)\in\text{graph of }T$. Hence\[
T(u+v)=Tu+Tv.
\]Pick any $a\in\mathbb{F}$, then $(au,aTu)\in\text{graph of }T$. Hence\[
T(au)=aTu.
\]Therefore, by definition, we conclude that $T$ is a linear map.
\end{proof}
\section{Exercise 3.E.2: $V_j$ is Finite-dimensional if $V_1\times\dots\times V_m$ is Finite-dimensional}
Suppose $V_1,\dots,V_m$ are vector spaces such that $V_1\times\dots\times V_m$ is finite-dimensional. Prove that $V_j$ is finite-dimensional for each $j=1,\dots,m$.
\begin{proof}
Suppose there exists $j\in\{1,\dots,m\}$ such that $V_j$ is infinite-dimensional. Then, by \ref{Exercise 2.A.14}, there is a sequence $v_{j1},v_{j2},\dots$ of vectors in $V_j$ such that $v_{j1},v_{j2},\dots,v_{jk}$ is linearly independent for every positive integer $k$.

Let the sequence $v_1,v_2,\dots$ of vectors in $V_1\times\dots\times V_m$ be defined by: for each positive integer $i$, the vector $v_i$ has its $j$\textsuperscript{th} slot equal to $v_{ji}$ and other slots equal to $0$. Then $v_{1},v_{2},\dots,v_{k}$ is linearly independent for every positive integer $k$. Therefore, again by \ref{Exercise 2.A.14}, $V_1\times\dots\times V_m$ is infinite-dimensional.
\end{proof}
\section{Exercise 3.E.3: Example of $U_1\times U_2$ Isomorphic to Non-Direct-Sum $U_1+U_2$}
Give an example of a vector space $V$ and subspaces $U_1,U_2$ of $V$ such that $U_1\times U_2$ is isomorphic to $U_1+U_2$ but $U_1+U_2$ is not a direct sum.
\begin{proof}
Let $V\equiv U_1\equiv\mathbb{R}^\infty$. Let\[
U_2\equiv\{(x,0,0,\dots):x\in\mathbb{R}\}.
\]Then $V$ is a vector space and $U_1,U_2$ are two subspaces of $V$, and\[
U_1+U_2=\mathbb{R}^\infty.
\]Since $U_1\cap U_2=U_2\ne\{0\}$, $U_1+U_2$ is not a direct sum. Let function $T: U_1\times U_2\to U_1+U_2$ be defined by\[
T((x,0,0,\dots),(y_1,y_2,\dots))=(x,y_1,y_2,\dots).
\]We claim that $T$ is a linear map. Suppose\[
u=((u_x,0,0,\dots),(u_{y_1},u_{y_2},\dots))
\]and\[
v=((v_x,0,0,\dots),(v_{y_1},v_{y_2},\dots)).
\]Then\[
\begin{split}
T(u+v)&=(u_x+v_x,u_{y_1}+v_{y_1},u_{y_2}+v_{y_2},\dots)\\
&=(u_x,u_{y_1},u_{y_2},\dots)+(v_x,v_{y_1},v_{y_2},\dots)\\
&=T(u)+T(v).
\end{split}
\]Pick any $a\in\mathbb{F}$,\[
\begin{split}
T(au)&=(au_x,au_{y_1},au_{y_2},\dots)\\
&=a(u_x,u_{y_1},u_{y_2},\dots)\\
&=aT(u).
\end{split}
\]Clearly $T$ is both injective and surjective. Therefore, $U_1\times U_2$ is isomorphic to $U_1+U_2$.
\end{proof}
\section{Exercise 3.E.4: $\mathcal{L}(V_1\times\dots\times V_m,W)$ and $\mathcal{L}(V_1, W)\times\dots\times\mathcal{L}(V_m,W)$ are Isomorphic}
Suppose $V_1,\dots,V_m$ are vector spaces. Prove that $\mathcal{L}(V_1\times\dots\times V_m,W)$ and $\mathcal{L}(V_1, W)\times\dots\times\mathcal{L}(V_m,W)$ are isomorphic vector spaces.
\begin{proof}
Let function $T:\mathcal{L}(V_1\times\dots\times V_m,W)\to\mathcal{L}(V_1, W)\times\dots\times\mathcal{L}(V_m,W)$ be defined by\[
T(P)=(P_1,\dots,P_m)
\]where\[
\begin{split}
P_1(v_1)&=P(v_1,0,\dots,0)\quad\forall v_1\in V_1,\\
P_2(v_2)&=P(0,v_2,\dots,0)\quad\forall v_2\in V_2,\\
&\vdots\\
P_m(v_m)&=P(0,0,\dots,v_m)\quad\forall v_m\in V_m.
\end{split}
\]

We claim that $T$ is a linear map. Pick any $P_1,P_2\in\mathcal{L}(V_1\times\dots\times V_m,W)$, then\[
\begin{split}
T(P_1+P_2)&=((P_1+P_2)_1,\dots,(P_1+P_2)_m)\\
&=((P_1)_1+(P_2)_1,\dots,(P_1)_m+(P_2)_m)\\
&=((P_1)_1,\dots,(P_1)_m) + ((P_2)_1,\dots,(P_2)_m)\\
&=T(P_1)+T(P_2).
\end{split}
\]Pick any $\lambda\in\mathbb{F}, P\in\mathcal{L}(V_1\times\dots\times V_m,W)$, then\[
\begin{split}
T(\lambda P)&=((\lambda P)_1,\dots,(\lambda P)_m)\\
&=\lambda(P_1,\dots,P_m)\\
&=\lambda T(P).
\end{split}
\]

We claim that $T$ is injective. Suppose $T(P)=0$. Then,\[
\begin{split}
P(v_1,0,\dots,0)&=0\quad\forall v_1\in V_1,\\
P(0,v_2,\dots,0)&=0\quad\forall v_2\in V_2,\\
&\vdots\\
P(0,0,\dots,v_m)&=0\quad\forall v_m\in V_m.
\end{split}
\]Hence, for all $v_1\in V_1, \dots, v_m\in V_m$, \[
P(v_1,\dots,v_m)=P(v_1,\dots,0)+\dots+P(0,\dots,v_m)=0.
\]Thus\[
P=0.
\]

We claim that $T$ is surjective. Pick any $(P_1,\dots,P_m)\in\mathcal{L}(V_1, W)\times\dots\times\mathcal{L}(V_m,W)$, let $P\in\mathcal{L}(V_1\times\dots\times V_m,W)$ be defined by\[
P(v_1,\dots,v_m)=P_1(v_1)+\dots+P_m(v_m)\quad\forall v_1\in V_1, \dots, v_m\in V_m.
\]Then\[
T(P)=(P_1,\dots,P_m).
\]

Therefore $T$ is an isomorphism from $\mathcal{L}(V_1\times\dots\times V_m,W)$ to $\mathcal{L}(V_1, W)\times\dots\times\mathcal{L}(V_m,W)$.
\end{proof}
\section{Exercise 3.E.5: $\mathcal{L}(V,W_1\times\dots\times W_m)$ and $\mathcal{L}(V, W_1)\times\dots\times\mathcal{L}(V,W_m)$ are Isomorphic}
Suppose $W_1,\dots,W_m$ are vector spaces. Prove that $\mathcal{L}(V,W_1\times\dots\times W_m)$ and $\mathcal{L}(V, W_1)\times\dots\times\mathcal{L}(V,W_m)$ are isomorphic vector spaces.
\begin{proof}
Let\[
\begin{split}
\Pi_1&\in\mathcal{L}(W_1\times\dots\times W_m, W_1),\\
\Pi_2&\in\mathcal{L}(W_1\times\dots\times W_m, W_2),\\
&\vdots\\
\Pi_m&\in\mathcal{L}(W_1\times\dots\times W_m, W_m)
\end{split}
\]be defined by\[
\begin{split}
\Pi_1(w_1,\dots,w_m)&=w_1,\\
\Pi_2(w_1,\dots,w_m)&=w_2,\\
&\vdots\\
\Pi_m(w_1,\dots,w_m)&=w_m
\end{split}
\]for all $(w_1,\dots,w_m)\in W_1\times\dots\times W_m$.

Let function $T:\mathcal{L}(V,W_1\times\dots\times W_m)\to\mathcal{L}(V, W_1)\times\dots\times\mathcal{L}(V,W_m)$ be defined by\[
T(P)=(P_{\Pi_1},\dots,P_{\Pi_m})
\]where\[
\begin{split}
P_{\Pi_1}(v)&=\Pi_1(P(v)),\\
P_{\Pi_2}(v)&=\Pi_2(P(v)),\\
&\vdots\\
P_{\Pi_m}(v)&=\Pi_m(P(v)).
\end{split}
\]for all $v\in V$.

We claim that $T$ is a linear map. Pick any $P_1,P_2\in\mathcal{L}(V,W_1\times\dots\times W_m)$, then\[
\begin{split}
T(P_1+P_2)&=((P_1+P_2)_{\Pi_1},\dots,(P_1+P_2)_{\Pi_m})\\
&=((P_1)_{\Pi_1}+(P_2)_{\Pi_1},\dots,(P_1)_{\Pi_m}+(P_2)_{\Pi_m})\\
&=((P_1)_{\Pi_1},\dots,(P_1)_{\Pi_m}) + ((P_2)_{\Pi_1},\dots,(P_2)_{\Pi_m})\\
&=T(P_1)+T(P_2).
\end{split}
\]Pick any $\lambda\in\mathbb{F}, P\in\mathcal{L}(V_1\times\dots\times V_m,W)$, then\[
\begin{split}
T(\lambda P)&=((\lambda P)_{\Pi_1},\dots,(\lambda P)_{\Pi_m})\\
&=\lambda(P_{\Pi_1},\dots,P_{\Pi_m})\\
&=\lambda T(P).
\end{split}
\]

We claim that $T$ is injective. Suppose $T(P)=0$. Then,\[
\begin{split}
\Pi_1(P(v))&=0,\\
\Pi_2(P(v))&=0,\\
&\vdots\\
\Pi_m(P(v))&=0.
\end{split}
\]for all $v\in V$. Hence, $P(v)=0$ for all $v\in V$. Thus\[
P=0.
\]

We claim that $T$ is surjective. Pick any $(P_1,\dots,P_m)\in\mathcal{L}(V, W_1)\times\dots\times\mathcal{L}(V,W_m)$, let $P\in\mathcal{L}(V,W_1\times\dots\times W_m)$ be defined by\[
P(v)=(P_1(v),\dots,P_m(v_m))\quad\forall v\in V.
\]Then\[
T(P)=(P_1,\dots,P_m).
\]

Therefore $T$ is an isomorphism from $\mathcal{L}(V,W_1\times\dots\times W_m)$ to $\mathcal{L}(V, W_1)\times\dots\times\mathcal{L}(V,W_m)$.
\end{proof}

\chapter{Duality}
\section{Exercise 3.F.1: Linear Functionals are either Surjective or Zero}
Explain why every linear functional is either surjective or the zero map.
\begin{proof}
Suppose $\varphi$ is a linear functional from $V$ to $\mathbb{F}$ and $\varphi$ is neither surjective nor the zero map. Then there exists $y_0\in\mathbb{F}$ such that\[
\varphi(v)\ne y_0,\quad\forall v\in V.
\]Also, there exist $x_0\in\mathbb{F}, v_0\in V$ such that\[
x_0\ne 0\text{ and }\varphi(v_0)=x_0.
\]

Since $\varphi$ is a linear map,\[
\begin{split}
\varphi(\frac{y_0}{x_0}\cdot v_0)&=\frac{y_0}{x_0}\cdot\varphi(v_0)\\
&=\frac{y_0}{x_0}\cdot x_0\\
&=y_0,
\end{split}
\]which contradicts with\[
\varphi(v)\ne y_0,\quad\forall v\in V.
\]
\end{proof}
\section{Exercise 3.F.2: Examples of Linear Functionals on $\mathbb{R}^{[0,1]}$}
Give three distinct examples of linear functionals on $\mathbb{R}^{[0,1]}$.
\begin{proof}
Pick any $c\in\mathbb{R}$. Let $\varphi:\mathbb{R}^{[0,1]}\to\mathbb{R}$ be defined by\[
\varphi(f)=f(c),\quad\forall f(x)\in\mathbb{R}^{[0,1]}.
\]Then we claim that $\varphi$ is a linear map. To show this, pick any $g(x), f(x)\in\mathbb{R}^{[0,1]}$. Then\[
\varphi(f+g)=f(c)+g(c)=\varphi(f)+\varphi(g).
\]Also, pick any $\lambda\in\mathbb{R}$. Then\[
\varphi(\lambda f)=(\lambda f)(c)=\lambda (f(c))=\lambda\varphi(f).
\]

Let $c\equiv1,2,3$ respectively, we then obtained three distinct linear functionals on $\mathbb{R}^{[0,1]}$.
\end{proof}
\section{Exercise 3.F.3: Existence of $\varphi\in V'$ such that $\varphi(v)=1$}
Suppose $V$ is finite-dimensional and $v\in V$ with $v\ne 0$. Prove that there exists $\varphi\in V'$ such that $\varphi(v)=1$.
\begin{proof}
Pick any basis $v_1,\dots,v_n$ of $V$ and its dual basis $\varphi_1,\dots,\varphi_n$ of $V'$, then there exists $a_1,\dots,a_n\in\mathbb{F}$ such that\[
v=a_1v_1+\dots+a_nv_n.
\]Since $v\ne 0$, there exists $k\in\{1,\dots,n\}$ such that\[
a_k\ne 0.
\]Let\[
\varphi\equiv\frac{1}{a_k}\varphi_k,
\]then $\varphi\in V'$ and \[
\begin{split}
\varphi(v)&=\frac{1}{a_k}\varphi_k(a_1v_1+\dots+a_nv_n)\\
&=\frac{1}{a_k}a_1\varphi_k(v_1)+\dots+\frac{1}{a_k}a_n\varphi_k(v_n)\\
&=\frac{1}{a_k}a_k\varphi_k(v_k)\\
&=\frac{1}{a_k}a_k\\
&=1.
\end{split}
\]
\end{proof}
\section{Exercise 3.F.10: The Dual Map Function that Takes $T$ to $T'$ is Linear}
Prove the first two bullet points in 3.101.
\begin{enumerate}
\item $(S+T)'=S'+T'$ for all $S,T\in\mathcal{L}(V,W)$.
\item $(\lambda T)'=\lambda T'$ for all $\lambda\in\mathbb{F}$ and all $T\in\mathcal{L}(V,W)$.
\end{enumerate}
\begin{proof}
\begin{enumerate}
\item Pick any $S,T\in\mathcal{L}(V,W),\varphi\in W',v\in V$. By definition, $(S+T)'(\varphi)=\varphi\circ(S+T)$. Since\[
\varphi\circ(S+T)(v)=\varphi(Sv+Tv)=\varphi(Sv)+\varphi(Tv),
\]we conclude that\[
\varphi\circ(S+T)=\varphi\circ S+\varphi\circ T=S'(\varphi)+T'(\varphi).
\]Therefore $(S+T)'=S'+T'$.
\item Pick any $\lambda\in\mathbb{F}$, any $T\in\mathcal{L}(V,W)$, any $\varphi\in W'$ and any $v\in V$. By definiton, $(\lambda T)'(\varphi)=\varphi\circ(\lambda T)$. Since\[
\varphi\circ(\lambda T)(v)=\varphi(\lambda Tv)=\lambda\varphi(Tv)=\lambda(\varphi\circ T)(v),
\]we conclude that\[
\varphi\circ(\lambda T)=\lambda(\varphi\circ T)=\lambda(T'(\varphi)).
\]Therefore $(\lambda T)'=\lambda T'$.
\end{enumerate}
\end{proof}
\section{Exercise 3.F.20: $U\subset W$ Implies $W^0\subset U^0$}
Suppose $U$ and $W$ are subsets of $V$ with $U\subset W$. Prove that $W^0\subset U^0$.
\begin{proof}
Pick any $\varphi\in W^0, w\in W, u\in U$. Since $U\subset W$, it follows that $u\in W$. Hence $\varphi(u)=0$ and $\varphi\in U^0$. Therefore $W^0\subset U^0$.
\end{proof}


\chapter{Polynomials}
\section{Exercise 4.A.2: $\{0\}\cup\{p\in\mathcal{P}(\mathbb{F}):\deg p=m\}$ is NOT a Subspace}
Suppose $m$ is a positive integer. Is the set\[
\{0\}\cup\{p\in\mathcal{P}(\mathbb{F}):\deg p=m\}
\]a subspace of $\mathcal{P}(\mathbb{F})$?
\begin{proof}
Denote\[
A_m\equiv \{0\}\cup\{p\in\mathcal{P}(\mathbb{F}):\deg p=m\}.
\]

Let $m\equiv 2$, and we claim that $A_2$ is not a subspace of $\mathcal{P}(\mathbb{F})$. To show this, let $p(z)\equiv z^2+z+1, q(z)\equiv -z^2+z+1$. Then both $p(z)$ and $q(z)$ are in $A_2$. However,\[
\deg \left(p(z)+q(z)\right)=\deg \left(2z+2\right)=1\ne 2.
\]Hence $p(z)+q(z)\notin A_2$.
\end{proof}
\section{Exercise 4.A.4: There Exists a $n$-Degree Polynomial with $m$ Zeros}
Suppose $m$ and $n$ are positive integers with $m\leqslant n$, and suppose $\lambda_1,\dots,\lambda_m\in\mathbb{F}$. Prove that there exists a polynomial $p\in\mathcal{P}(F)$ with $\deg p=n$ such that $0=p(\lambda_1)=\dots=p(\lambda_m)$ and such that $p$ has no other zeros.
\begin{proof}
Let\[
p(z)\equiv(z-\lambda_1)\dots (z-\lambda_m)^{n-m+1}.
\]Then, by definition, $\deg p=m-1+(n-m+1)=n$ and $0=p(\lambda_1)=\dots=p(\lambda_m)$. Hence it only suffices to show that $p$ has no other zeros.

Suppose $p$ has another zero $\tau$ such that\[
\tau\ne\lambda_i,\quad\forall i\in\{1,\dots,m\}.
\]Then\[
0\ne\tau-\lambda_i,\quad\forall i\in\{1,\dots,m\}
\]and\[
0\ne(\tau-\lambda_1)\dots (\tau-\lambda_m)^{n-m+1}.
\]However,\[
p(\tau)=(\tau-\lambda_1)\dots (\tau-\lambda_m)^{n-m+1},
\]which contradicts with\[
p(\tau)=0.
\]
\end{proof}
		
\backmatter
\bibliographystyle{amsalpha}
\bibliography{minqi_s_solutions_to_the_exercises_of_sheldon_axler_2014_linear_algebra_done_right_3rd_edition}
\end{document}
